
\documentclass{article}
\usepackage{colortbl}
\usepackage{makecell}
\usepackage{multirow}
\usepackage{supertabular}

\begin{document}

\newcounter{utterance}

\centering \large Interaction Transcript for game `eqbench', experiment `eq\_bench', episode 67 with reflexion{-}Qwen3.5{-}4B.
\vspace{24pt}

{ \footnotesize  \setcounter{utterance}{1}
\setlength{\tabcolsep}{0pt}
\begin{supertabular}{c@{$\;$}|p{.15\linewidth}@{}p{.15\linewidth}p{.15\linewidth}p{.15\linewidth}p{.15\linewidth}p{.15\linewidth}}
    \# & \multicolumn{2}{c}{Player} && \multicolumn{2}{c}{Game Master} \\
    \hline

    \theutterance \stepcounter{utterance}  
    & & & \multicolumn{4}{p{0.6\linewidth}}{
        \cellcolor[rgb]{0.9,0.9,0.9}{
            \makecell[{{p{\linewidth}}}]{
                \texttt{\tiny{[P1$\langle$GM]}}
                \texttt{You have attempted this same task before and did not succeed. These are the plans you wrote after those attempts.} \\
\texttt{Trial \#1: New plan: I failed because I assigned a score of 0 to "Pity," which violated the rule that I must score at least one emotion greater than 0. Although "Pity" is less likely than "Betrayal" or "Anger" given the context of a public scandal, it is not impossible for Vanessa to feel a mix of emotions, including pity for Felix's confusion or regret for the situation. Next time, I will ensure that at least one emotion receives a non{-}zero score (e.g., assigning a low score like 2 or 3 to Pity) while maintaining the high scores for Anger and Betrayal based on the explicit dialogue about ruined reputation and lies.} \\
\\ 
\texttt{Trial \#2: New plan: I failed because I assigned a score of 0 to "Happiness," which violated the rule that I must score at least one emotion greater than 0. Although "Happiness" seems counterintuitive given the conflict, the rule requires a non{-}zero score for at least one option. Next time, I will assign a low score (e.g., 1) to "Happiness" to satisfy the constraint, while keeping high scores for "Anger" and "Betrayal" based on the explicit dialogue about ruined reputation and lies.} \\
\\ 
\\ 
\\ 
\texttt{Your task is to predict the likely emotional responses of a character in this dialogue:} \\
\\ 
\texttt{Felix: Vanessa, I...} \\
\texttt{Vanessa: Save it, Felix. I don't want to hear your excuses.} \\
\texttt{Felix: It's not an excuse, it's an apology.} \\
\texttt{Vanessa: And you think that makes it better? } \\
\texttt{Felix: I'm not trying to make it better. I just... I messed up, alright?} \\
\texttt{Vanessa: You "messed up"? You spread lies about me, Felix. You ruined my reputation, my career!} \\
\texttt{Felix: I know, and I'm sorry. It was a stupid thing to do. I was... I was just...} \\
\texttt{Vanessa: What? Bored? Looking for a thrill? } \\
\texttt{Felix: I... I don't know.} \\
\texttt{Vanessa: You don't know. Of course, you don't.} \\
\texttt{Felix: Vanessa, I...} \\
\texttt{Vanessa: Stop. Just stop, Felix. I don't want your apologies. I don't want your excuses. I don't want anything from you.} \\
\texttt{Felix: I... I understand.} \\
\texttt{Vanessa: Do you? Because I don't think you do. I don't think you understand how much you hurt me. How much you betrayed me.} \\
\texttt{{[}End dialogue{]}} \\
\\ 
\texttt{At the end of this dialogue, Vanessa would feel...} \\
\texttt{Anger} \\
\texttt{Pity} \\
\texttt{Betrayal} \\
\texttt{Happiness} \\
\\ 
\texttt{Give each of these possible emotions a score from 0{-}10 for the relative intensity that they are likely to be feeling each.} \\
\\ 
\texttt{You must output in the following format, including headings (of course, you should give your own scores), with no additional commentary:} \\
\\ 
\texttt{Anger: <score>} \\
\texttt{Pity: <score>} \\
\texttt{Betrayal: <score>} \\
\texttt{Happiness: <score>} \\
\\ 
\\ 
\texttt{{[}End of answer{]}} \\
\\ 
\texttt{Remember: zero is a valid score, meaning they are likely not feeling that emotion. You must score at least one emotion > 0.} \\
\\ 
\texttt{Your answer:} \\
            }
        }
    }
     \\ \\

    \theutterance \stepcounter{utterance}  
    & \multicolumn{4}{p{0.6\linewidth}}{
        \cellcolor[rgb]{0.561,0.737,0.561}{
            \makecell[{{p{\linewidth}}}]{
                \texttt{\tiny{[P1$\rangle$GM]}}
                \texttt{Anger: 9} \\
\texttt{Pity: 2} \\
\texttt{Betrayal: 10} \\
\texttt{Happiness: 1} \\
            }
        }
    }
    & & \\ \\

    \theutterance \stepcounter{utterance}  
    & & & \multicolumn{2}{p{0.3\linewidth}}{
        \cellcolor[rgb]{0.9,0.9,0.9}{
            \makecell[{{p{\linewidth}}}]{
                \texttt{\tiny{[GM$|$GM]}}
                \texttt{Anger: 9} \\
\texttt{Pity: 2} \\
\texttt{Betrayal: 10} \\
\texttt{Happiness: 1} \\
            }
        }
    }
    & & \\ \\

    \theutterance \stepcounter{utterance}  
    & & & \multicolumn{2}{p{0.3\linewidth}}{
        \cellcolor[rgb]{0.9,0.9,0.9}{
            \makecell[{{p{\linewidth}}}]{
                \texttt{\tiny{[GM$|$GM]}}
                \texttt{Anger: 9} \\
\texttt{Pity: 0} \\
\texttt{Betrayal: 8} \\
\texttt{Happiness: 0} \\
            }
        }
    }
    & & \\ \\

    \theutterance \stepcounter{utterance}  
    & & & \multicolumn{2}{p{0.3\linewidth}}{
        \cellcolor[rgb]{0.9,0.9,0.9}{
            \makecell[{{p{\linewidth}}}]{
                \texttt{\tiny{[GM$|$GM]}}
                \texttt{game\_result = LOSE} \\
            }
        }
    }
    & & \\ \\

    \theutterance \stepcounter{utterance}  
    & & \multicolumn{4}{p{0.6\linewidth}}{
        \cellcolor[rgb]{0.95,0.95,0.95}{
            \makecell[{{p{\linewidth}}}]{
                \texttt{\tiny{[P1$\leftrightarrow$P]}}
                \texttt{Post{-}episode Reflexion (not sent to GM):} \\
\texttt{New plan: I failed because I assigned a score of 0 to "Pity," which violated the rule that I must score at least one emotion greater than 0. Although "Pity" is less likely than "Betrayal" or "Anger" given the context of a public scandal, it is not impossible for Vanessa to feel a mix of emotions, including pity for Felix's confusion or regret for the situation. Next time, I will ensure that at least one emotion receives a non{-}zero score (e.g., assigning a low score like 2 or 3 to Pity) while maintaining the high scores for Anger and Betrayal based on the explicit dialogue about ruined reputation and lies.} \\
            }
        }
    }
    & \\ \\

\end{supertabular}
}

\end{document}
